\documentclass[aspectratio=169,11pt]{beamer}

\usetheme{Madrid}
\usecolortheme{default}
\usefonttheme{professionalfonts}
\setbeamertemplate{navigation symbols}{}
\setbeamertemplate{footline}[frame number]

\usepackage{amsmath, amssymb, booktabs, array}

\title{Nonstandard Preferences in Labor Supply, Effort, Savings, and Training}
\subtitle{Behavioral Labor --- Week 2}
\author{Teaching deck draft}
\date{}

\begin{document}

\begin{frame}
  \titlepage
\end{frame}

\begin{frame}{Central question}
\begin{itemize}
    \item How do nonstandard preferences alter labor supply, effort, savings, and training?
    \item Week 1 gave the full Behavioral Labor map.
    \item Week 2 isolates the preference branch before beliefs, expectations, attention, firms, and policy design.
    \item The discipline remains: preference object, labor margin, identifying variation, welfare.
\end{itemize}
\end{frame}

\begin{frame}{Bridge from Week 1}
\small
\begin{tabular}{p{0.30\linewidth} p{0.58\linewidth}}
\toprule
Week 1 object & Week 2 narrowing \\
\midrule
Benchmark labor choice & Workers choose hours, effort, saving, training, and task completion \\
Behavioral taxonomy & Focus only on nonstandard preferences \\
Identification habit & Name the margin and the counterfactual \\
Welfare caution & Observed choices may not reveal stable welfare rankings \\
\bottomrule
\end{tabular}
\end{frame}

\begin{frame}{DellaVigna taxonomy: Week 2 highlighted}
\small
\begin{tabular}{p{0.29\linewidth} p{0.58\linewidth}}
\toprule
Branch & Labor interpretation \\
\midrule
\textbf{Nonstandard preferences} & \textbf{Present bias, commitment demand, risk aversion, reference dependence, reciprocity, fairness} \\
Nonstandard beliefs & Subjective expectations, perceived returns, optimism, job-search beliefs \\
Nonstandard decision-making & Attention, salience, complexity, limited consideration, defaults \\
Responses & Firms, markets, and policies adapt to predictable behavioral patterns \\
\bottomrule
\end{tabular}

\vspace{0.5em}
\textbf{This week:} preferences mapped to labor supply timing, effort intensity, savings, and training completion.
\end{frame}

\begin{frame}{Preference choice object}
\[
a_i^B \in \arg\max_{a \in \mathcal{A}_i}
\tilde{U}_i(a;\beta_i,\rho_i,r_i,\eta_i)
\]

\small
\begin{tabular}{p{0.16\linewidth} p{0.70\linewidth}}
\toprule
Parameter & Interpretation \\
\midrule
$\beta_i$ & present bias and self-control \\
$\rho_i$ & risk exposure or curvature relevant for saving, pay, and training \\
$r_i$ & reference point for earnings, effort, targets, or evaluations \\
$\eta_i$ & social preferences, reciprocity, fairness, employer surplus \\
\bottomrule
\end{tabular}
\end{frame}

\begin{frame}{Present bias}
\[
U_t = u(c_t,\ell_t) + \beta \sum_{\tau=t+1}^{T} \delta^{\tau-t}u(c_\tau,\ell_\tau),
\qquad \beta \in (0,1]
\]

\begin{itemize}
    \item Preference object: current utility receives extra weight relative to future utility.
    \item Labor margin: effort today, hours timing, training completion, saving now.
    \item Empirical design: compare planned or committed future behavior to realized behavior when current costs arrive.
    \item Key caution: distinguish self-control from liquidity constraints, fatigue, and learning.
\end{itemize}
\end{frame}

\begin{frame}{Commitment demand at work}
\small
\begin{tabular}{p{0.26\linewidth} p{0.28\linewidth} p{0.30\linewidth}}
\toprule
Object & Operationalization & Observed margin \\
\midrule
Self-control problem & Choice of dominated or restrictive work contract & Output and contract take-up \\
Payday timing & Effort around paydays & Daily or weekly productivity \\
Training commitment & Deadlines, deposits, cohorts & Attendance and completion \\
Savings commitment & Payroll deduction or default enrollment & Participation and contribution rates \\
\bottomrule
\end{tabular}

\vspace{0.6em}
Kaur, Kremer, and Mullainathan make commitment a labor-design object: the question is not only whether incentives work, but whether workers value discipline against future selves.
\end{frame}

\begin{frame}{Savings and self-control in work-linked environments}
\begin{itemize}
    \item Savings is labor-relevant when tied to employment, payroll systems, retirement, or worker planning.
    \item Default enrollment designs observe participation and contribution behavior.
    \item Savings-incentive designs observe take-up and contribution response when returns become salient.
    \item Interpretation problem: low saving can reflect procrastination, transaction costs, information, liquidity, or true low demand.
\end{itemize}
\end{frame}

\begin{frame}{Risk aversion: keep it worker-side}
\small
\begin{tabular}{p{0.30\linewidth} p{0.30\linewidth} p{0.27\linewidth}}
\toprule
Labor object & Preference role & Design burden \\
\midrule
Retirement saving & Precaution and portfolio risk & Separate risk tastes from information and default effects \\
Variable pay & Exposure to output or earnings risk & Separate risk aversion from effort costs and selection \\
Training & Uncertain returns to self-investment & Separate risk from credit constraints and expected returns \\
\bottomrule
\end{tabular}

\vspace{0.5em}
Risk aversion matters here, but it should not turn Week 2 into an occupational-sorting lecture.
\end{frame}

\begin{frame}{Reference dependence and loss aversion}
\[
u(x \mid r) = m(x) + \mu\!\left(m(x)-m(r)\right)
\]

\begin{itemize}
    \item Preference object: outcomes are evaluated relative to a reference point.
    \item Labor margin: target labor supply, performance after wage disappointment, training milestones.
    \item Identifying variation: earnings targets, arbitration shocks, wage expectations, target changes.
    \item Caution: rule out standard nonlinear incentives, income effects, fatigue, learning, and composition.
\end{itemize}
\end{frame}

\begin{frame}{Social preferences and workplace effort}
\[
U_i = w_i - c(e_i) + \eta_i S(e_i,\pi_f,g_f)
\]

\begin{itemize}
    \item Preference object: worker utility includes fairness, reciprocity, employer surplus, or social meaning.
    \item Labor margin: effort, productivity, extra work, cooperation.
    \item Identifying variation: gifts, employer returns, whether work benefits the employer, pay framing.
    \item Caution: distinguish social preferences from monitoring, repeated-game incentives, and selection.
\end{itemize}
\end{frame}

\begin{frame}{Empirical designs and what they identify}
\scriptsize
\begin{tabular}{p{0.23\linewidth} p{0.24\linewidth} p{0.22\linewidth} p{0.20\linewidth}}
\toprule
Preference & Variation & Margin & Main threat \\
\midrule
Present bias & contract choice, payday timing, deadlines & effort, completion & liquidity, fatigue \\
Commitment & restrictive contracts, defaults, deposits & take-up, future action & selection, monitoring \\
Savings self-control & default enrollment, match incentives & participation, contribution & information, transaction costs \\
Reference dependence & targets, arbitration, wage surprises & hours, performance & standard incentives \\
Social preferences & gifts, employer surplus, pay framing & effort, productivity & career concerns, supervision \\
\bottomrule
\end{tabular}
\end{frame}

\begin{frame}{Welfare and policy interpretation}
\begin{itemize}
    \item Present bias: ex ante and ex post selves can disagree.
    \item Commitment: restrictions can raise effort while reducing future flexibility.
    \item Defaults: participation may rise without proving every worker is better off.
    \item Reference dependence: welfare depends on how reference points are formed or manipulated.
    \item Social preferences: higher effort can be meaningful motivation, costly pressure, or surplus transfer.
\end{itemize}
\end{frame}

\begin{frame}{Week 2 lab}
\begin{itemize}
    \item Reproduce: synthetic Kaur-Kremer-Mullainathan-style workplace self-control factbook.
    \item Diagnose: contract choice, output, payday timing, and alternative explanations.
    \item Transfer: synthetic work-linked savings default design.
    \item Challenge: Mas-style pay reference points and performance.
\end{itemize}
\end{frame}

\begin{frame}{Bridge to Week 3}
\begin{itemize}
    \item Week 2 treats behavior as a preference problem.
    \item Week 3 asks when the observed behavior is instead a beliefs or expectations problem.
    \item Core next margins: job search, perceived returns, wage expectations, and training payoffs.
    \item Same empirical habit: object, margin, variation, counterfactual, welfare.
\end{itemize}
\end{frame}

\begin{frame}{Takeaways}
\begin{itemize}
    \item Nonstandard preferences become labor economics only when mapped to concrete worker-side margins.
    \item Present bias, commitment, reference dependence, and social preferences imply different empirical designs.
    \item Savings and training are labor-facing when connected to work, payroll, retirement, and self-investment.
    \item Welfare requires more than an observed behavioral response.
\end{itemize}
\end{frame}

\end{document}
